\begin{papercard}{mayring2025}
\cf{Stages} \stagetag{L} (a sceptical field test of ChatGPT as a qualitative-content-analysis tool), \stagetag{C} (what rule-guided qualitative content analysis requires: category definition, level of abstraction, units), \stagetag{I} (re-run consistency as intra-/inter-coder agreement), \stagetag{H} (data-protection risk of sending transcripts to an LLM). A 13-page \emph{FQS} field report (Jan 2025, 34 numbered paragraphs) by the developer of the qualitative content analysis method that the thesis may cite. It is an anecdotal test by one author on one short German interview excerpt, translated to English, so its value is as an expert's argument and as a counter-weight to the enthusiastic quantitative annotation papers.

\cf{Aim and design} The author tests ChatGPT versions 3.5 and 4 on a qualitative content analysis of an interview excerpt, using a naive request, precise prompts, a re-run, and a larger set of interviews \pc{mayring2025}{1}{I describe a series of test runs to examine the contribution that the AI-based program ChatGPT in both Versions 3.5. and 4 can make to a qualitative content analysis of interview texts}. Motive: a lecturer advised students to paste interviews into ChatGPT and ask for a content analysis ``according to Mayring''. The test text is an interview with a former East-German teacher made unemployed after reunification; the \emph{sample solution} is produced by the author's own method (inductive category formation) \pc{mayring2025}{3}{These rules result in six inductive categories}.

\cf{The reference method} Mayring's qualitative content analysis assigns categories ``as precisely as possible, which can be verified by coding agreements'', with clear procedural steps for each of its techniques (summarisation, inductive category formation, explication, deductive category assignment, mixed forms), so that a model solution exists \pc{mayring2025}{3}{the aim is to assign categories as precisely as possible, which can be verified by coding agreements}. Sample task: ``What stress factors due to unemployment were reported?'' Technique: inductive category formation with the category definition ``Subjectively negatively experienced aspects of changes in the life situation due to unemployment'' \pc{mayring2025}{3}{Subjectively negatively experienced aspects of changes in the life situation due to unemployment} and an abstraction level ``as concrete as possible in the text, but not so specific that the category could only be coded for the person''. Six categories K1--K6: arbitrary dismissal, loss of the teaching profession, household and garden not fulfilling, unprepared for new conditions after the fall of communism, uncertainty about the lawsuit, financial uncertainty. Segments were left uncoded where the point was merely a reply to an interviewer question (suggestive) or too interpretative for the abstraction level.

\cf{Test 1: naive request} Prompt (translated): ``Please do a qualitative content analysis according to Mayring with the following interview excerpt'' \pc{mayring2025}{4}{Please do a qualitative content analysis according to Mayring with the following interview excerpt}; the author calls this ``a completely inadequate, naïve request, but one that is likely to occur frequently''. ChatGPT 3.5 (24 March 2024) gave a four-step procedure (transcription, categorisation, coding, interpretation) that ``is incorrect, has not been published anywhere in my work'' and resembles the approach of Schreier that Mayring criticises \pc{mayring2025}{5}{The presentation of the method with the four steps listed is incorrect, has not been published anywhere in my work}. Of the two negative categories, four of the six reference categories were roughly found with three miscategorisations \pc{mayring2025}{5}{four of the six categories from the sample solution are roughly found in my model solution, a rather poor result, with three miscategorizations}.

\cf{Test 2: other concepts} Asked for Schreier's, Kuckartz's and Braun \& Clarke's approaches, the tool returned structurally similar outputs; for Schreier it produced open/axial/selective coding from grounded theory, described incorrectly; for Braun \& Clarke the answer claimed to find hidden patterns but stayed descriptive; asked for ``Brown and Harris'' (a depression-research pair, not text analysts) instead of Braun and Clarke, it produced a ``thematic analysis'' with meaningless results in both versions \pc{mayring2025}{6}{Chat GPT nevertheless carried out an analysis in both Versions 3.5 and 4.0 and presented more or less meaningless results}. Irony, indirect allusions and a slip of the tongue were not recognised \pc{mayring2025}{6}{ironic remarks and indirect allusions were not recognized}; complex output requests (frequency lists) were ignored; about three transcript pages was the limit for the free version \pc{mayring2025}{6}{I quickly reached a limit of around three transcript pages}.

\cf{Test 3: precise prompt} Prompt with the exact research question, the inductive technique, the abstraction rule and the ``subjectively significant'' filter. Output: five categories, two of them not independent factors, three reference categories missed \pc{mayring2025}{7}{The answer from ChatGPT 3.5 (from March 24, 2024) was significantly worse compared to the sample solution}. A popular ``Mayring Mentor'' custom GPT contained serious errors (only one technique, wrong use of paraphrasing, only the coding unit defined, main-category formation skipped, no references) \pc{mayring2025}{7}{The template contained serious errors: Only one of the techniques we developed was described, namely inductive category formation}. The author notes that iterating on prompts requires method expertise, and that if one has precise ideas to write good prompts one could analyse the material directly without AI \pc{mayring2025}{5}{If you already had precise ideas at the beginning in order to be able to send meaningful prompts, you could carry out the analysis directly on the material without AI}.

\cf{Test 4: re-run (coder agreement)} Identical task a day later (25 March 2024) gave ten categories, adding non-stress items such as ``use of support services'' and changing others; ``Overall, the coding agreement test must be considered completely failed''; also ChatGPT 4 gave different solutions for identical input \pc{mayring2025}{8}{Overall, the coding agreement test must be considered completely failed}. The author is unsure whether a re-run measures intra-coder (same person twice) or inter-coder agreement \pc{mayring2025}{8}{I am not quite sure whether a re-run of the identical input in ChatGPT would rather measure intracoder agreement}; both matter as quality criteria.

\cf{Test 5: ChatGPT 4, about 10 pages (six interviews)} Request for stress factors across six cases with frequencies (18 March 2024): 20 categories, each ``1 time'', contrary to the instruction to count frequencies \pc{mayring2025}{9}{A list of 20 inductive categories was provided, but contrary to the instructions, each was coded only once}. Compared with the sample solution (16 categories): five similar, but frequencies wrong; four of six interviewees considered the dismissal unfair and the tool missed this \pc{mayring2025}{10}{Four of the six interviewees reported that they considered their dismissal to be unfair, which ChatGPT 4 did not recognize}; clear stress passages not coded; categories listed with no clear text passage. Capacity: 4,096 tokens (about three pages) in the free version, about 500 pages in paid version 4, but recommended in portions, whereas qualitative content analyses often cover thousands of pages.

\cf{Conclusion} Results are ``rough approximations'' in the best case and contain ``gross errors''; closest to the sample solution was ChatGPT 3.5 with short text and simple questions; re-analysis with identical questions yielded different results; the potential is far from exhausted but a precise-prompt approach needs expertise; iteration does not match the rule-guided style of the method \pc{mayring2025}{10}{An iterative approach to the results also does not correspond to the research style of the qualitative content analysis developed by our working group}. Data protection: sending an interview to ChatGPT requires the interviewee's explicit written consent, and GDPR compliance is ``extremely questionable'' even for a 30-year-old consented interview \pc{mayring2025}{10}{sending an interview transcript to ChatGPT actually requires the explicit, written consent of the interviewee}. What it can do: rough content categorisation and topic overview, suitable for descriptive, thematic approaches (Kuckartz, Braun \& Clarke) \pc{mayring2025}{11}{rough content categorizations of text passages are reasonably successful and that an overview of topics in the material can be given}; less suitable for rule-guided techniques, for theory-based in-depth analyses (reflexive TA, grounded theory, psychoanalytic interpretation); editing an LLM first pass doesn't save work because errors are unknown, and a rough result can distort the analysis \pc{mayring2025}{11}{you don't know exactly where the inaccuracies and errors are, so you have to work through all the material in any case}. Tools change rapidly, so this is an ``interim review'' \pc{mayring2025}{12}{tools such as ChatGPT are constantly undergoing revisions, which is why this can only be an interim review}.

\cf{Limits (card)} One short German excerpt translated into English, a handful of prompts, one rater (the author, whose own method is the reference and who has an interest in it), no systematic metric (no $\kappa$, F1, precision/recall), chat interface instead of API with fixed parameters (so temperature and version are unknown), mostly free version 3.5, two dates in March 2024, no use of few-shot examples or a full codebook as in Pangakis or Törnberg, and a task (open inductive category formation) different from a closed deductive coding task with a fixed codebook; so it cannot be generalised to deductive coding of screens or policies. Its arguments about method fit, units of analysis and ethics are nevertheless independent of the model version.

\cf{Use for this thesis} (1) A qualified expert voice that the thesis should cite when stating that an LLM is not used for inductive, interpretive coding, and that LLM output is not treated as a qualitative content analysis. (2) Differentiates the thesis' use: deductive coding with a closed codebook (stage C/I/L) rather than open category development; here the relevant questions are the ones Mayring tests (does the output follow the coding rules? is a re-run consistent? are units and abstraction level respected?). (3) Re-run instability as an observation to be measured with temperature, seed and repeated runs (cf.\ Rodriguez, Pangakis). (4) Ethics: do not send interview material or participant data to a commercial LLM without consent and a lawful basis (relevant if the thesis includes participant material; public store and policy text is less sensitive). (5) A reminder to name the exact procedure (e.g.\ ``deductive category assignment'') when the thesis cites Mayring, since the paper shows how often LLMs mix up named methods.
\end{papercard}
